% !TeX program = pdflatex
% Document autonome : les résultats numériques changent avec les données,
% mais les définitions et les formules décrites ici restent reproductibles.
%
% Ce fichier est la spécification de référence des calculs publiés sur
% https://radar-parlementaire.fr. Toute mesure affichée par le site doit être
% définie ici ; l'inverse n'est pas exigé. Quand un calcul change dans
% src/radar/, ce document change dans le même commit.
\documentclass[11pt,a4paper]{article}

\usepackage[T1]{fontenc}
\usepackage[utf8]{inputenc}
\usepackage[french]{babel}
\usepackage{lmodern}
\usepackage[a4paper,margin=2.45cm]{geometry}
\usepackage{amsmath,amssymb,mathtools}
\usepackage{booktabs,array}
\usepackage{xcolor}
\usepackage{fancyhdr}
\PassOptionsToPackage{hyphens}{url}  % coupure des URL longues (avant hyperref)
\usepackage{microtype}
\usepackage{hyperref}   % chargé en dernier, comme le veut son manuel
\Urlmuskip=0mu plus 1mu\relax

\hypersetup{
  colorlinks=true,
  linkcolor=blue!55!black,
  urlcolor=blue!55!black,
  citecolor=blue!55!black,
  pdftitle={Mesurer l'activité de l'Assemblée nationale à partir de ses
            données ouvertes : définitions, dénominateurs et incertitudes},
  pdfauthor={Sylvain Desroziers},
  pdfsubject={Note méthodologique du Radar parlementaire},
  pdfkeywords={Assemblée nationale, données ouvertes, scrutins publics,
               points idéaux, cohésion de groupe, cosignatures, transparence,
               reproductibilité},
  pdflang={fr-FR},
  bookmarksnumbered=true,
  % La table des matières imprimée s'arrête aux sections : à raison de quatre
  % sous-sections chacune, la détailler coûtait trois pages sur dix-neuf. Les
  % signets du lecteur PDF, eux, gardent les deux niveaux — c'est là qu'on
  % navigue vraiment.
  bookmarksdepth=2,
}
\setcounter{tocdepth}{1}

% ── mise en page d'article ────────────────────────────────────────────────
\pagestyle{fancy}
\fancyhf{}
\fancyhead[L]{\footnotesize\itshape Radar parlementaire — note méthodologique}
\fancyhead[R]{\footnotesize\thepage}
\renewcommand{\headrulewidth}{0.4pt}
\fancypagestyle{plain}{%
  \fancyhf{}\renewcommand{\headrulewidth}{0pt}\fancyfoot[C]{\footnotesize\thepage}}

\setlength{\parskip}{0.55em}
\setlength{\parindent}{0pt}
\renewcommand{\arraystretch}{1.22}
\newcommand{\ind}{\mathbf{1}}
\newcommand{\Pp}{\mathbb{P}}
\newcommand{\E}{\mathbb{E}}
\newcommand{\R}{\mathbb{R}}
\newcommand{\code}[1]{\texttt{#1}}
\newcommand{\important}[1]{\begin{quote}\small\color{blue!55!black}\textbf{À retenir.} #1\end{quote}}
\newcommand{\prudence}[1]{\begin{quote}\small\color{red!55!black}\textbf{Prudence d'interprétation.} #1\end{quote}}

\title{\vspace{-1.2cm}
Mesurer l'activité de l'Assemblée nationale\\à partir de ses données ouvertes\\[0.35em]
\large Définitions, dénominateurs et incertitudes des indicateurs du
\emph{Radar parlementaire}}

\author{%
  \normalsize\textbf{Sylvain Desroziers}\\[0.15em]
  \small Citoyen. Sans affiliation institutionnelle, sans financement, sans parti.\\
  \small\href{mailto:sylvain.desroziers@protonmail.com}
             {sylvain.desroziers@protonmail.com}\\[0.5em]
  \small\itshape Rédaction assistée par Claude (Anthropic) ;
  voir la déclaration en section~\ref{sec:ia}.
}
\date{\small\today\\[0.3em]
  \small Version de référence : \href{https://radar-parlementaire.fr/methode.html}
                                     {radar-parlementaire.fr/methode.html}}

\begin{document}
\maketitle
\thispagestyle{plain}

\begin{abstract}
\noindent
L'Assemblée nationale publie en données ouvertes le détail nominatif de ses
scrutins publics, ses amendements et la composition de ses groupes. Ces
fichiers sont accessibles à tous et lus par presque personne : leur format
suppose des compétences que la citoyenneté ne devrait pas exiger. Le
\emph{Radar parlementaire} les transforme en indicateurs consultables, et ce
document spécifie ces indicateurs.

Nous décrivons chaque mesure publiée par le site : sa définition formelle, son
numérateur, son dénominateur, la population sur laquelle elle porte, les
observations qu'elle exclut et l'incertitude qui l'entoure. Sont couverts
l'accord de vote entre députés et entre groupes, la participation, la cohésion
et la dissidence, deux réductions descriptives (analyse en composantes
principales et positionnement multidimensionnel), un modèle de points idéaux
estimé avec ses intervalles par bootstrap de blocs, le traitement de
l'abstention, le réseau de cosignatures d'amendements et la détection de
sujets sous contrôle du taux de fausses découvertes.

Deux principes gouvernent l'ensemble. Aucun taux n'est publié sans le
dénombrement qui le produit : un pourcentage privé de son dénominateur n'est
pas une information, c'est une opinion présentée comme un fait. Et aucun
indicateur ne délivre de verdict : ces mesures décrivent une part observable
de l'activité parlementaire, elles n'évaluent ni la qualité d'un élu ni la
justesse d'une position. Le document est écrit pour être vérifié, y compris —
et surtout — par qui souhaite démontrer qu'il se trompe.
\end{abstract}

\vspace{0.4em}
\noindent\small\textbf{Mots-clés} — données ouvertes parlementaires ; scrutins
publics ; points idéaux ; cohésion de groupe ; cosignatures ; transparence
algorithmique ; reproductibilité.
\normalsize

\vspace{0.6em}
\hrule
\vspace{0.4em}

\tableofcontents
\newpage

\section{Introduction}

\subsection{Une transparence qui existe sans être accessible}

L'\emph{open data} de l'Assemblée nationale est une réussite institutionnelle
réelle et une impasse pratique. Chaque scrutin public y figure avec le vote
nominatif de chaque député ; chaque amendement avec son auteur, ses
cosignataires et son sort. Rien n'est caché. Mais la donnée est distribuée en
archives compressées de plusieurs centaines de mégaoctets, en JSON issus d'une
conversion XML dont les conventions ne sont documentées nulle part, et sans
aucun indicateur dérivé. La transparence est juridiquement acquise et
matériellement hors de portée.

L'espace laissé vide se remplit tout seul. Il se remplit de chiffres cités
sans source, de classements de « députés les plus absents » dont le
dénominateur n'est jamais donné, de pourcentages d'accord calculés sur des
échantillons choisis après coup. Ces chiffres circulent parce qu'ils sont les
seuls disponibles sous une forme lisible, et non parce qu'ils sont justes.

Le \emph{Radar parlementaire} occupe cet espace autrement : il recalcule
chaque indicateur depuis les fichiers officiels et publie, avec chaque
résultat, ce qui permet de le contester.

\subsection{Ce que ce document ajoute au site}

Un site qui affiche « 42~\% de présence » demande à être cru. Un site qui
affiche « 42~\%, soit 178 scrutins sur 424 où l'élu pouvait voter, après
retrait des scrutins antérieurs à son entrée en fonction et des non-votes
imposés par une fonction incompatible » demande à être vérifié. La seconde
formulation ne vaut que si la règle qu'elle invoque est écrite quelque part,
précisément, une fois pour toutes. C'est ce document.

Il remplit trois fonctions distinctes.

\textbf{Une spécification.} Chaque mesure affichée par le site est définie ici
formellement. Le code implémente cette spécification ; en cas de divergence,
c'est le code qui a tort. Cette contrainte a déjà servi : plusieurs
incohérences entre deux chemins de calcul d'une même grandeur ont été
découvertes en écrivant les définitions plutôt qu'en lisant le code.

\textbf{Un support pédagogique.} Les notions mobilisées — dénominateur
conditionnel, indice de cohésion, variable latente, bootstrap de blocs,
correction pour tests multiples — ne sont pas hors de portée d'un lecteur
attentif, mais elles s'expliquent mal en légende d'un graphique. Elles sont
introduites ici avec des exemples de petite taille, calculables à la main.

\textbf{Une prise pour la contradiction.} C'est la fonction la plus
importante. Un indicateur sur la vie politique sera contesté, et il doit
l'être : un député dont le taux de dissidence est publié a le droit de savoir
comment il a été obtenu, sur quels scrutins, avec quelle définition de la
ligne de son groupe, et à partir de quel seuil le chiffre cesse d'être
significatif. Sans ce document, cette contestation ne peut porter que sur la
conclusion. Avec lui, elle porte sur la méthode — c'est-à-dire qu'elle devient
un débat qu'on peut trancher.

\subsection{Trois exigences}

\textbf{Reproductibilité.} Les données sources sont publiques, le code est
publié sous licence libre, et les tables intermédiaires sont archivées à
chaque mise à jour. Un lecteur disposant d'une machine ordinaire peut
reconstruire l'intégralité des chiffres du site.

\textbf{Réfutabilité.} Chaque indicateur est accompagné de sa population de
référence, de ses exclusions et, lorsque cela a un sens, de son incertitude.
Une mesure dont l'échantillon est trop mince n'est pas affichée avec une
précaution en petits caractères : elle n'est pas affichée.

\textbf{Absence de verdict.} Le projet ne classe pas les députés en bons et en
mauvais. Il ne pondère pas les scrutins par leur importance supposée, il ne
construit pas de note de synthèse, il n'infère aucune intention. La section
\ref{sec:lire} récapitule, indicateur par indicateur, ce que chacun ne
démontre pas. Cette retenue n'est pas de la tiédeur : c'est la condition pour
qu'un chiffre reste utilisable par des lecteurs qui ne partagent pas les mêmes
convictions.

\subsection{Ce dont ce travail hérite}

Aucune des méthodes employées ici n'est nouvelle, et c'est délibéré : sur un
objet politique, l'originalité méthodologique est un risque, pas un atout.
L'analyse quantitative des scrutins nominatifs est un domaine constitué depuis
près d'un siècle. L'indice de cohésion d'un groupe remonte à
Rice~\cite{rice1928}. L'estimation d'une position latente à partir de votes
binaires suit la lignée de Poole et Rosenthal~\cite{poole1985,poole1997}, dans
la variante probabiliste formalisée par Clinton, Jackman et
Rivers~\cite{clinton2004} ; l'application à un parlement européen
multipartisan, plus proche du cas français que le bipartisme américain, doit
beaucoup à Hix, Noury et Roland~\cite{hix2006}. Les outils statistiques
auxiliaires sont classiques : indice de similarité de Jaccard~\cite{jaccard1901},
positionnement multidimensionnel de Torgerson~\cite{torgerson1952}, bootstrap
par blocs de Künsch~\cite{kunsch1989} pour tenir compte de la dépendance entre
scrutins d'une même séance, et contrôle du taux de fausses découvertes de
Benjamini et Hochberg~\cite{benjamini1995} pour la détection de sujets.

L'apport de ce projet n'est donc pas méthodologique. Il est dans l'application
systématique de ces méthodes à la législature en cours, dans leur mise à jour
automatique, et dans la publication intégrale des dénominateurs.

\subsection{Plan}

La section~\ref{sec:question} pose les unités d'observation et le périmètre.
La section~\ref{sec:cube} décrit le passage des fichiers bruts aux matrices de
vote. Les sections~\ref{sec:proximite} à~\ref{sec:sujets} définissent les
indicateurs, dans l'ordre où le site les présente. La
section~\ref{sec:lire} récapitule les lectures abusives à éviter, et la
section~\ref{sec:reproduire} indique comment tout reconstruire.

\section{Avant toute formule : quelle est la question ?}\label{sec:question}

Un calcul politique n'est jamais seulement une formule. Il suppose aussi une
question, une unité d'observation et un périmètre. Ici, les unités élémentaires
sont :

\begin{itemize}
  \item un \textbf{scrutin public} $j$ ;
  \item un \textbf{député} $i$ ;
  \item un \textbf{vote nominatif} de $i$ au scrutin $j$ ;
  \item lorsque les amendements sont disponibles, un \textbf{amendement} et
  l'ensemble de ses signataires.
\end{itemize}

Le projet répond à des questions descriptives : « qui vote de la même façon ?
», « un groupe a-t-il une ligne nette ? », « quels termes deviennent plus
présents ? », « qui cosigne avec qui ? ». Il ne permet pas, à lui seul, de
répondre à « pourquoi ce député a-t-il voté ainsi ? », « ce député travaille-t-il
bien ? » ou « deux groupes sont-ils alliés ? ». Ces dernières questions
nécessitent du contexte politique, qualitatif et institutionnel.

\important{Un pourcentage n'est interprétable qu'avec son dénominateur et sa
population de référence. C'est pourquoi les fonctions du projet renvoient les
comptes bruts avec les taux.}

\subsection{Qui fait quoi : les fonctions qui déplacent un chiffre}
\label{sec:fonctions}

Plusieurs calculs de ce document dépendent de la fonction exercée par un député.
Ces fonctions sont familières au monde parlementaire et opaques pour tout le
monde sinon ; les voici, avec l'endroit exact où chacune intervient. Ce n'est
pas une présentation de l'Assemblée : c'est la liste des rôles dont l'oubli
fausse un résultat.

\begin{center}
\begin{tabular}{>{\raggedright\arraybackslash}p{.22\linewidth}p{.40\linewidth}p{.30\linewidth}}
\toprule
Fonction & Ce que c'est & Où elle change un chiffre \\
\midrule
\textbf{Président de l'Assemblée} & Élu par ses pairs pour la durée de la législature, il dirige l'institution. Son siège surélevé a donné le mot \emph{perchoir}, employé comme synonyme de la fonction. Il ne prend pas part aux votes qu'il préside. & Non-vote structurel \code{PAN} : retiré du dénominateur de présence (§\ref{sec:participation}). \\
\addlinespace
\textbf{Président de séance} & La présidence tourne au cours de la journée : des vice-présidents se relaient au perchoir. Celui qui préside à l'instant du scrutin ne vote pas. & Non-vote structurel \code{PSE}, même retrait. \\
\addlinespace
\textbf{Membre du Gouvernement} & Un député nommé ministre cesse de siéger et son suppléant le remplace ; la source continue néanmoins d'enregistrer ses non-votes. & Non-vote structurel \code{MG}, même retrait. \\
\addlinespace
\textbf{Rapporteur} & Député chargé d'examiner un texte au nom d'une commission et d'en proposer les corrections. Deux co-rapporteurs cosignent par dizaines des amendements purement rédactionnels, quels que soient leurs groupes. & Amendements de rapporteur exclus du réseau de cosignatures (§\ref{sec:cosign}). \\
\addlinespace
\textbf{Commission} & Formation restreinte où les textes sont examinés avant la séance publique. L'essentiel du travail législatif s'y fait, le plus souvent \textbf{sans vote nominatif publié}. & Hors périmètre : invisible dans tous les indicateurs. \\
\addlinespace
\textbf{Groupe politique} & Réunion d'au moins quinze députés. Il structure le temps de parole, les postes et les positions de vote. & Population de référence de la cohésion, de la ligne et de la dissidence. \\
\bottomrule
\end{tabular}
\end{center}

Les trois premières lignes ont un point commun qui justifie leur traitement
identique : dans les trois cas, \textbf{c'est l'institution qui empêche le vote,
pas l'élu qui s'y dérobe}, et la source le déclare. La cinquième ligne dit
l'inverse : le travail en commission est réel, massif, et ce document n'en
mesure rien — non par choix, mais parce qu'il ne laisse pas de trace nominative
exploitable.

\prudence{La conséquence pratique est qu'un député très investi en commission,
ou occupé à présider, peut apparaître peu actif ici. Le taux de présence répond
à « a-t-il voté quand il pouvait voter ? », jamais à « travaille-t-il ? ».}

\section{Des fichiers bruts au cube de votes}\label{sec:cube}

\subsection{Nettoyer sans inventer de données}

Les fichiers de l'Assemblée sont des JSON issus d'un format XML. Ils peuvent
encoder une liste d'un seul élément comme un dictionnaire, ou une valeur absente
comme un objet particulier. Le module \code{src/radar/parse.py} homogénéise ces
formes puis produit des tables Parquet : \code{deputes}, \code{scrutins},
\code{votes}, \code{positions\_groupe} et, si demandé, \code{amendements}.

La table \code{votes} a une ligne par couple observé $(i,j)$ et quatre états
possibles : \emph{pour}, \emph{contre}, \emph{abstention} et \emph{non-votant}.
Une absence de ligne dans une table n'est jamais convertie silencieusement en
un vote contre ou en une absence : ce serait fabriquer une observation.

\subsection{Une matrice par position}

Pour $n$ députés et $m$ scrutins sélectionnés, le projet construit quatre
matrices booléennes de taille $n\times m$ :
\[
 P_{ij}=\ind(\text{$i$ vote pour au scrutin $j$}),\qquad
 C_{ij}=\ind(\text{$i$ vote contre au scrutin $j$}),
\]
\[
 A_{ij}=\ind(\text{$i$ s'abstient au scrutin $j$}),\qquad
 N_{ij}=\ind(\text{$i$ est \emph{non-votant} au scrutin $j$}).
\]
Le vote est dit \emph{exprimé} lorsqu'il vaut pour, contre ou abstention :
\[
 V_{ij}=P_{ij}+C_{ij}+A_{ij}.
\]
Ces quatre matrices sont exclusives : un couple $(i,j)$ relève d'au plus une
d'entre elles. $V_{ij}$ vaut donc $1$ si le député a exprimé une position, et
$0$ sinon.

\subsection{Le non-votant : la quatrième matrice, et pourquoi elle compte}
\label{sec:nonvotant}

$N_{ij}$ demande plus d'explications que les trois autres, car elle recouvre
une situation que le vocabulaire courant confond avec l'absence.

\textbf{Cinq états, et non quatre.} Un couple $(i,j)$ peut se trouver dans cinq
situations, dont la cinquième n'a pas de matrice parce qu'elle se déduit des
autres :

\begin{center}
\begin{tabular}{>{\raggedright\arraybackslash}p{.20\linewidth}p{.36\linewidth}p{.34\linewidth}}
\toprule
État & Ce que contient la source & Codage \\
\midrule
Pour / contre / abstention & une ligne avec la position & $P$, $C$ ou $A$ vaut 1 \\
Non-votant & une ligne, position \code{nonVotant}, \textbf{plus un code de cause} & $N_{ij}=1$ \\
Absent & \textbf{aucune ligne} & les quatre matrices valent 0 \\
\bottomrule
\end{tabular}
\end{center}

La différence entre les deux dernières lignes est la clé. Un \emph{absent} est
un député dont la source ne dit rien : il pouvait voter, il ne l'a pas fait, et
on ignore pourquoi. Un \emph{non-votant} est un député que l'Assemblée déclare
comme n'ayant pas pris part au vote, \textbf{en indiquant le motif}. Ce motif
n'est pas une explication a posteriori : c'est une donnée du fichier source.

\textbf{Les motifs observés.} Sur la 17\textsuperscript{e} législature, la
source n'emploie que trois codes, et tous trois désignent une fonction
incompatible avec le fait de voter :

\begin{center}
\begin{tabular}{>{\raggedright\arraybackslash}p{.10\linewidth}p{.46\linewidth}p{.34\linewidth}}
\toprule
Code & Fonction & Portée \\
\midrule
\code{MG}  & membre du Gouvernement & un ministre ne vote pas \\
\code{PSE} & président de séance & celui qui préside au moment du scrutin \\
\code{PAN} & président de l'Assemblée nationale & le titulaire du perchoir \\
\bottomrule
\end{tabular}
\end{center}

On note $\mathcal S$ l'ensemble de ces codes et
$\mathrm{cause}_{ij}$ le motif déclaré. Un non-vote est dit \textbf{structurel}
lorsque $\mathrm{cause}_{ij}\in\mathcal S$.

\textbf{Pourquoi $N_{ij}$ n'entre pas dans $V_{ij}$.} Un non-votant n'a exprimé
aucune position : il ne peut donc être ni d'accord ni en désaccord avec
quiconque, et l'inclure dans un taux d'accord fabriquerait une concordance là
où il n'y a pas de vote. $N$ n'apparaît pour cette raison ni dans $V$, ni dans
le numérateur d'accord $K$ de la section~\ref{sec:proximite}.

\textbf{Mais $N_{ij}$ n'est pas pour autant sans usage} : c'est elle qui permet
de distinguer, au dénominateur de présence, ce qu'un député n'a pas fait de ce
qu'on l'a empêché de faire. Elle réapparaît sous la forme du terme $R_i$ de la
section~\ref{sec:participation}, et c'est son seul emploi.

\important{Un non-vote structurel n'est pas une absence : c'est l'institution
qui interdit le vote, pas l'élu qui s'y soustrait. Le confondre avec de
l'absentéisme revient à reprocher à un député d'avoir présidé la séance ou
d'être entré au Gouvernement.}

\prudence{La réciproque est tout aussi importante : le retrait ne vaut que pour
ce que la source justifie explicitement. Élargir $\mathcal S$ au-delà des motifs
déclarés permettrait de faire disparaître n'importe quelle absence en la
requalifiant, et le dénominateur cesserait d'être opposable.}

\subsection{La matrice signée, et ce qu'elle écrase}
\label{sec:signee}

Pour les analyses factorielles exploratoires, le projet utilise enfin une
matrice signée :
\[
 S_{ij}=P_{ij}-C_{ij}\in\{-1,0,+1\}.
\]
Cette matrice a une propriété qu'il faut énoncer avant de s'en servir : elle
\textbf{n'a que trois valeurs pour cinq états}. Le $0$ y recouvre à lui seul
l'abstention, le non-vote structurel, l'absence, et le cas d'un député non
encore élu. Ces quatre situations sont politiquement très différentes et
deviennent indiscernables.

Ce n'est pas une remarque de principe. Sur la 17\textsuperscript{e}
législature, un scrutin public réunit en moyenne moins du tiers de
l'Assemblée : \textbf{environ trois quarts des cellules de $S$ valent $0$ sans
que personne se soit abstenu}. Autrement dit, $S$ décrit d'abord qui était là,
et seulement ensuite qui a voté quoi. Les conséquences sur l'ACP sont détaillées
en section~\ref{sec:acp} ; elles sont assez sérieuses pour que la matrice signée
soit réservée aux cartes descriptives et jamais employée à situer un député en
particulier.

\subsection{L'éligibilité : ne pas pénaliser un mandat qui n'avait pas commencé}
\label{sec:eligibilite}

Le cube contient enfin $E_{ij}$, l'indicatrice que le député $i$ \emph{occupait
son siège} le jour du scrutin $j$. Cette matrice est essentielle pour les
mesures de présence : un élu arrivé en cours de législature ne pouvait pas
participer aux scrutins antérieurs, et le lui reprocher serait le principal
biais possible du site.

\paragraph{Un mandat est une suite de périodes, et non un intervalle.}
Un député nommé au Gouvernement quitte son siège ; son suppléant l'occupe ;
puis le titulaire le reprend. La source publie alors \emph{deux} mandats
séparés par un trou de plusieurs mois. On note donc
$\mathcal{P}_i=\{[a_{i1},b_{i1}],\dots,[a_{ik_i},b_{ik_i}]\}$ les périodes
d'occupation du siège par $i$, et
\[
 E_{ij}=\ind\!\left(\exists\,[a,b]\in\mathcal{P}_i \;:\;
 a\leq d_j\leq b\right),
\]
une \emph{union} d'intervalles. Pour une période encore ouverte, $b$ est
considérée comme future. Les bornes sont closes des deux côtés : la source date
la sortie du titulaire au jour $b$ et l'entrée de son suppléant au jour $b+1$,
de sorte que deux périodes voisines ne se recouvrent jamais d'une journée.

\paragraph{Les bornes viennent de la prise de fonction, jamais de
\code{dateDebut}.} Le champ \code{dateDebut} de l'open data ne date pas
l'entrée en fonction : l'Assemblée y inscrit l'ouverture de la législature pour
tous les mandats, y compris ceux de remplaçants entrés deux ans plus tard --
584 des 677 mandats de la 17\textsuperscript{e} législature portent la même
valeur. La date d'entrée effective est \code{mandature/datePriseFonction}, et
c'est elle qui borne $\mathcal{P}_i$.

\paragraph{Ces deux points ont été publiés à tort.} Le site a longtemps
reconstruit $E_{ij}$ comme un intervalle unique allant de la première entrée à
la dernière sortie, assis sur \code{dateDebut}. Deux conséquences, toutes deux
visibles : le trou d'un mandat interrompu était comblé, et un remplaçant était
réputé siéger depuis le premier jour de la législature. 261\,214 couples
député\,$\times$\,scrutin étaient déclarés éligibles à tort, 86 députés
recevaient des absences qu'ils ne pouvaient pas commettre, et une page de
scrutin de juin 2026 annonçait 581 députés en fonction là où l'Assemblée en
compte 577. Le correctif et son ampleur figurent au registre des corrections du
site.

\paragraph{Un invariant, désormais vérifié avant publication.} L'article 24 de
la Constitution plafonne l'effectif de l'Assemblée à 577 députés. Pour toute
date de scrutin,
\[
 \sum_i E_{ij}\leq 577,
\]
et deux députés d'une même circonscription ne peuvent avoir de période commune.
Ces deux contraintes, plus la concordance de nos totaux avec ceux publiés par
l'Assemblée, sont contrôlées sur les 8\,434 scrutins avant la génération du
site : un échec arrête la publication au lieu de l'assortir d'un avertissement.
C'est le module \code{radar.controles}, et il ne vérifie pas des formules mais
des effectifs -- aucun test de calcul n'aurait pu voir l'erreur ci-dessus, car
un effectif faux ne contredit aucune formule.

\subsection{Le parcours d'un texte, en quatre mots}
\label{sec:parcours}

La suite parle d'amendements, d'articles et de « votes sur l'ensemble ». Ces
termes désignent des moments distincts de la fabrication d'une loi, et c'est
leur ordre qui explique pourquoi ils ne pèsent pas pareil.

\begin{enumerate}
  \item Un \textbf{texte} est déposé : \emph{projet de loi} s'il vient du
  Gouvernement, \emph{proposition de loi} s'il vient d'un parlementaire. Il est
  découpé en \textbf{articles}, numérotés.
  \item Chaque article peut être modifié par des \textbf{amendements} — des
  propositions de retouche, parfois d'un seul mot. N'importe quel député peut en
  déposer, et un texte disputé en attire des milliers.
  \item L'Assemblée vote, dans l'ordre : d'abord les amendements un par un, puis
  chaque article, puis enfin \textbf{l'ensemble du texte}. Ce dernier vote est le
  seul qui engage sur la loi entière.
  \item Le texte part au Sénat, revient éventuellement — c'est la
  \emph{navette} —, et en cas de désaccord persistant une \emph{commission mixte
  paritaire} réunit les deux chambres pour chercher un compromis.
\end{enumerate}

Cette mécanique produit une asymétrie massive et purement arithmétique :
\textbf{une seule loi très disputée génère des milliers de votes d'amendement
et un unique vote sur l'ensemble}. Sur la 17\textsuperscript{e} législature,
environ six scrutins publics sur sept portent sur un amendement.

Une conséquence suit immédiatement. Un vote d'amendement est souvent tactique :
on peut approuver un aménagement technique tout en combattant le texte, ou
rejeter une retouche que l'on juge insuffisante sans s'opposer au principe.
Agréger tous les scrutins « toutes portées confondues » revient donc à mesurer
surtout de la tactique parlementaire, et à laisser une loi très amendée peser
plus lourd que deux ans de votes solennels.

\subsection{Portée politique d'un scrutin}

C'est pourquoi le projet classe les scrutins par \textbf{portée}, selon ce sur
quoi ils portent. Le titre est classé par motifs réguliers en une catégorie,
puis en une portée :

\begin{center}
\begin{tabular}{>{\raggedright\arraybackslash}p{.25\linewidth}p{.30\linewidth}p{.32\linewidth}}
\toprule
Catégorie & Portée & Lecture \\
\midrule
Amendement, sous-amendement & \code{detail} & Travail technique et tactique sur le texte. \\
Article, motion de procédure, déclaration & \code{intermediaire} & Enjeu intermédiaire. \\
Ensemble d'un texte, motion de censure & \code{texte} & Vote qui engage le plus nettement une position sur le texte. \\
\bottomrule
\end{tabular}
\end{center}

Le classement est déterministe : il dépend de la syntaxe du titre, pas d'un
codage politique fait à la main. Les analyses peuvent être recalculées séparément
pour les trois portées. Cette séparation évite qu'une longue séquence
d'amendements écrase, par son seul nombre, les votes sur l'ensemble d'un texte.

\subsection{Ordinaire, solennel, motion de censure : le type est autre chose}

La portée décrit \emph{sur quoi} l'Assemblée vote. Le \textbf{type} décrit
\emph{comment} elle vote, et la source le publie séparément. Les deux ne se
recouvrent pas, et il faut les tenir distincts pour ne pas compter deux fois la
même précaution.

\begin{center}
\begin{tabular}{>{\raggedright\arraybackslash}p{.10\linewidth}p{.30\linewidth}p{.50\linewidth}}
\toprule
Code & Type & Ce qui le caractérise \\
\midrule
\code{SPO} & scrutin public ordinaire & Le cas courant. Le vote a lieu depuis le pupitre, au fil de la séance, souvent devant un hémicycle clairsemé. \\
\code{SPS} & scrutin public \textbf{solennel} & Annoncé à l'avance, inscrit à l'ordre du jour, généralement le mardi après les questions au Gouvernement. La présence y est bien plus forte. \\
\code{MOC} & motion de censure & Procédure distincte : seuls les soutiens de la motion votent, et l'adoption exige la majorité absolue des membres de l'Assemblée, non des votants. \\
\bottomrule
\end{tabular}
\end{center}

Le scrutin solennel est rare — quelques dizaines sur une législature — et
presque toujours de portée \code{texte} : on solennise le vote sur l'ensemble
d'une loi, pas sur un sous-amendement. L'inverse n'est pas vrai : la majorité
des votes sur l'ensemble d'un texte reste des scrutins ordinaires.

Ce type est disponible comme filtre (\code{types\_vote}), mais les indicateurs
publiés ne l'utilisent pas : la portée est le découpage retenu, parce qu'elle
sépare des populations bien plus grandes et repose sur le contenu du vote
plutôt que sur son protocole.

\prudence{La motion de censure ne se lit pas comme les autres scrutins. Ne pas
la voter n'est pas s'y opposer : c'est le mode normal d'expression de ceux qui
ne la soutiennent pas. La source n'y publie d'ailleurs ni « contre » ni
abstention, et un taux de participation calculé dessus n'aurait pas de sens.}

Le projet calcule aussi la \emph{contestation} du scrutin :
\[
 q_j=\frac{\min(n_{\mathrm{pour},j},n_{\mathrm{contre},j})}
 {\max(n_{\mathrm{pour},j}+n_{\mathrm{contre},j},1)}.
\]
Elle est proche de $0$ pour un scrutin quasi unanime et de $0{,}5$ pour un
scrutin très partagé. On peut l'utiliser pour retirer les votes joués d'avance,
qui gonflent mécaniquement l'accord général.

\section{Proximité de vote : compter les accords, puis seulement les diviser}\label{sec:proximite}

\subsection{La définition entre deux députés}

Pour une paire de députés $(i,k)$, le nombre de scrutins où les deux ont
exprimé un suffrage est
\[
 M_{ik}=\sum_{j=1}^{m}V_{ij}V_{kj}.
\]
Le nombre d'accords est
\[
 K_{ik}=\sum_{j=1}^{m}
 \left(P_{ij}P_{kj}+C_{ij}C_{kj}+A_{ij}A_{kj}\right).
\]
Le taux d'accord est donc
\[
 T_{ik}=\frac{K_{ik}}{M_{ik}}.
\]
Une absence ou un non-vote ne compte ni comme accord ni comme désaccord : il
sort du dénominateur. L'abstention, en revanche, est une position et deux
abstentions concordent.

Dans le code, les matrices de tous les $K_{ik}$ et $M_{ik}$ sont obtenues d'un
coup par produits matriciels :
\[
 K=PP^{\mathsf T}+CC^{\mathsf T}+AA^{\mathsf T},\qquad M=VV^{\mathsf T}.
\]
C'est exactement la formule précédente, calculée pour toutes les paires à la
fois.

\subsection{Exemple jouet}

Considérons six scrutins. A et B votent toujours pareil ; C est opposé à A,
sauf qu'ils s'abstiennent tous les deux au cinquième scrutin ; D ne vote qu'une
fois, comme A.

\begin{center}
\begin{tabular}{c*{6}{c}}
\toprule
 & $s_1$ & $s_2$ & $s_3$ & $s_4$ & $s_5$ & $s_6$ \\
\midrule
A & pour & pour & contre & contre & abst. & pour \\
B & pour & pour & contre & contre & abst. & pour \\
C & contre & contre & pour & pour & abst. & contre \\
D & pour & -- & -- & -- & -- & -- \\
\bottomrule
\end{tabular}
\end{center}

Alors $T_{AB}=6/6=100\%$, $T_{AC}=1/6\approx16{,}7\%$ et
$T_{AD}=1/1=100\%$. Le dernier résultat est exact mais peu informatif. Le
projet masque un taux quand $M_{ik}$ est inférieur à un seuil, fixé par défaut
à 30 scrutins communs (et à 20 pour les seuls votes de portée \code{texte}).
La diagonale est également masquée : un député est toujours d'accord avec lui-même.

\important{Un « 100\% d'accord » sans le nombre de scrutins communs est souvent
un faux signal. Le seuil ne rend pas une paire peu observée moins proche : il
reconnaît que sa proximité n'est pas mesurée de façon suffisamment stable.}

\subsection{De deux députés à deux groupes : deux questions légitimes}

Soient $G$ et $H$ deux groupes. Après exclusion des paires sous le seuil, le
projet publie deux conventions :
\begin{align*}
 \overline T_{GH}&=\frac{1}{|\mathcal P_{GH}|}
 \sum_{(i,k)\in\mathcal P_{GH}}T_{ik},\\
 T^{\text{agrégé}}_{GH}&=
 \frac{\sum_{(i,k)\in\mathcal P_{GH}}K_{ik}}
 {\sum_{(i,k)\in\mathcal P_{GH}}M_{ik}}.
\end{align*}
Ici $\mathcal P_{GH}$ est l'ensemble des paires retenues, avec une paire interne
à $G$ comptée une seule fois.

La première formule est une \textbf{moyenne non pondérée des paires}. Elle
répond à : « si je tire un député de $G$ et un de $H$, quelle est leur proximité
moyenne ? ». La seconde est un \textbf{quotient de sommes} : une paire observée
sur 400 scrutins y pèse quatre fois une paire observée sur 100. Elle répond à :
« sur tous les votes communs de toutes les paires, quelle part est concordante ? ».

Ce ne sont pas deux estimations concurrentes d'une vérité unique ; elles ont des
pondérations et donc des questions différentes. Le nombre de paires et le
nombre de scrutins communs accompagnent les deux mesures.

\subsection{Cohésion interne}

La cohésion d'un groupe est la convention $\overline T_{GG}$ : l'accord moyen
entre les paires distinctes de ses membres. Elle mesure une homogénéité de votes,
pas l'existence d'une consigne ni l'intensité des débats internes. Un effectif
doit toujours être affiché à côté : une moyenne calculée sur trois binômes est
beaucoup plus instable qu'une moyenne calculée sur plusieurs milliers.

\subsection{Comparer les portées et vérifier la taille d'échantillon}

Pour une paire de groupes, le projet recalcule les taux sur les trois ensembles
disjoints de scrutins : détail, intermédiaire et texte. Un écart entre la portée
\code{detail} et la portée \code{texte} peut cependant venir de deux causes :
un vrai changement de comportement ou le fait que le second ensemble est plus
petit.

La contre-épreuve est un test de rééchantillonnage. Si la portée \code{texte}
contient $m_T$ scrutins, on tire sans remise $m_T$ scrutins parmi les scrutins de
détail, $B$ fois (par défaut $B=40$), et l'on recalcule l'accord. Si le taux
observé sur les textes est extérieur au minimum--maximum des $B$ tirages, il est
peu plausible que la seule taille d'échantillon explique l'écart. Cette procédure
est un contrôle intuitif ; elle ne remplace pas un modèle causal.

\section{Présence, ligne de groupe et dissidence}

\subsection{Participation : un numérateur, deux retraits au dénominateur}
\label{sec:participation}

Pour le député $i$, le nombre de votes exprimés est
\[
 X_i=\sum_j V_{ij}E_{ij}.
\]
Le nombre de scrutins éligibles est $L_i=\sum_jE_{ij}$. Le nombre de non-votes
structurels reprend la matrice $N$ et l'ensemble de motifs $\mathcal S$
introduits en section~\ref{sec:nonvotant} :
\[
 R_i=\sum_j N_{ij}\,E_{ij}\,
 \ind\!\left(\mathrm{cause}_{ij}\in\mathcal S\right).
\]
Le dénominateur et le taux sont alors
\[
 D_i=\max(L_i-R_i,0),
 \qquad \mathrm{participation}_i=\frac{X_i}{D_i},
\]
avec une valeur manquante si $D_i=0$.

Les deux retraits répondent à deux objections distinctes. Le premier ($E_{ij}$)
écarte les scrutins antérieurs — ou postérieurs — au mandat, \emph{et ceux tenus
pendant une interruption} : on ne reproche pas à un élu arrivé en cours de
législature les votes tenus avant son arrivée, ni à un ancien ministre ceux
tenus pendant que son suppléant occupait son siège (section~
\ref{sec:eligibilite}). Le second ($R_i$) écarte les non-votes que
l'institution impose. Le plancher à zéro
n'a pas de rôle statistique : un non-vote structurel suppose un mandat en cours,
donc $R_i\leq L_i$ par construction, et le \code{max} ne protège que d'une
source incohérente.

\textbf{Un cas limite qui n'en est pas un.} Le retrait $R_i$ peut être du même
ordre que $L_i$, et c'est attendu. Le président de l'Assemblée ne vote pas tant
qu'il préside : sur les votes qui engagent de la 17\textsuperscript{e}
législature, la présidente est déclarée non-votante pour cause \code{PAN} sur la
quasi-totalité des scrutins, et a exprimé un suffrage à chacun des scrutins
restants. Son taux est donc de 100~\% sur un dénominateur de quelques dizaines
de votes, et non de 8~\% sur l'ensemble. Les deux nombres sont issus des mêmes
données ; un seul répond à la question « à quelle fraction des scrutins où cet
élu pouvait voter a-t-il voté ? ».

\important{Ce cas illustre pourquoi le dénominateur doit être publié à côté du
taux, et non seulement calculé. « 100~\% sur 19 votes » et « 8~\% sur 221 »
décrivent la même personne ; sans les effectifs, le lecteur ne peut pas savoir
lequel des deux répond à sa question.}

\prudence{Un $D_i$ petit rend le taux instable : à 19 scrutins, un vote de plus
ou de moins déplace le taux de cinq points. Le dénominateur sert à lire le taux,
pas seulement à le justifier.}

\subsubsection*{Publier un taux et le comparer sont deux choses}

De cette instabilité découle une règle qui n'est pas une exclusion. Un taux dont
le dénominateur est inférieur à un seuil $D_{\min}$ — fixé à 20, le même nombre
qu'exige déjà un taux d'accord sur les votes qui engagent — reste calculé,
reste publié, et reste accompagné de son effectif. Ce qu'il ne fait plus, c'est
entrer dans la \emph{distribution de comparaison} : la médiane, le maximum, les
rangs et la représentation graphique de la population.

La raison est arithmétique. Ces quantités servent de repère à tous les autres
députés, et un taux mesuré sur quelques scrutins les fixe sans les mériter. Le
cas s'est produit : la présidente de l'Assemblée, qui ne vote pas tant qu'elle
préside, a exprimé un suffrage à chacun des dix-neuf scrutins où elle n'était
pas empêchée. Son taux de 100~\% est exact, et il devenait « le député le plus
assidu de la législature », au-dessus de qui a voté 190 fois sur 245. Le maximum
de l'Assemblée était fixé par son plus petit dénominateur.

\important{Écarter un taux d'une comparaison n'est pas le juger faible. C'est
reconnaître qu'il ne répond pas à la même question que les autres, faute de
porter sur la même assiette. La fiche concernée l'écrit explicitement plutôt que
d'afficher un rang que rien ne fonde.}

Le taux mesure la participation aux scrutins publics, et rien d'autre : ni la
présence en commission, ni le travail en circonscription, ni les interventions
en séance sans vote.

\subsection{Le vote par délégation : ce que « présence » recouvre}
\label{sec:delegation}

Un député empêché peut confier son vote à un collègue, qui l'exprime en son
nom ; le règlement en autorise une par député. Le vote est alors juridiquement
celui du délégant, et il entre au numérateur $X_i$ de la
section~\ref{sec:participation} exactement comme un vote émis en personne.

Cette convention est la seule défendable — la délégation transfère l'exercice
du vote, pas son imputation — mais elle a une conséquence que le mot
« présence » masque. On publie donc le compte séparément. En notant
$\Delta_{ij}$ l'indicatrice d'un vote délégué :
\[
 \delta_i=\sum_j V_{ij}E_{ij}\Delta_{ij},
 \qquad
 \text{part déléguée}_i=\frac{\delta_i}{X_i}.
\]
Le dénominateur est ici $X_i$, le nombre de suffrages exprimés, et non $D_i$ :
la question posée est « parmi les votes émis en son nom, combien l'ont été par
un tiers ? », qui porte sur le numérateur de la présence et non sur son
assiette.

\textbf{Les ordres de grandeur importent pour lire ce chiffre.} La délégation
n'est pas un cas limite : elle représente environ 15~\% des suffrages exprimés
sur l'ensemble des scrutins, et près de 23~\% sur les seuls votes qui engagent
— c'est-à-dire précisément l'assiette des mesures principales de ce projet. La
médiane par député y avoisine 20~\%, et les valeurs extrêmes dépassent 90~\%.

\important{Le taux de la section~\ref{sec:participation} mesure les suffrages
\emph{émis au nom} du député, non sa présence physique dans l'hémicycle. Pour
un député dont la moitié des votes sont délégués, ces deux lectures divergent
franchement, et c'est la première qui est exacte.}

\textbf{Ce que la source ne permet pas.} Le fichier de scrutin nomme le
délégant et signale que son vote a été délégué ; il ne nomme jamais le collègue
qui l'a porté. On ne peut donc ni classer les récipiendaires, ni reconstituer
les binômes de délégation. Le site affiche pour cette raison un compte et une
part, jamais un porteur : la donnée n'existe pas, et la déduire serait
l'inventer.

\prudence{Une part déléguée élevée ne se lit pas comme un manquement. Elle peut
résulter d'une mission, d'un congé, d'une maladie ou d'une fonction — autant de
motifs que la source ne donne pas. Le chiffre dit comment le vote a été émis,
jamais pourquoi.}

\subsection{Reconstruire la ligne à partir du dépouillement}

Pour un groupe $g$ et un scrutin $j$, on recompte les trois effectifs :
\[
 n^+_{gj},\quad n^-_{gj},\quad n^0_{gj},\qquad
 n_{gj}=n^+_{gj}+n^-_{gj}+n^0_{gj}.
\]
La part de la position dominante est
\[
 h_{gj}=\frac{\max(n^+_{gj},n^-_{gj},n^0_{gj})}{n_{gj}}.
\]
La source publie une position majoritaire de groupe, mais le projet ne l'utilise
pas pour classer les députés : il recalcule les comptes à partir des votes
nominatifs. La ligne observée est la position dominante uniquement si :

\begin{enumerate}
  \item il y a au moins cinq votants dans le groupe ;
  \item il n'y a pas d'égalité au sommet ;
  \item $h_{gj}>0{,}5$.
\end{enumerate}

Pourquoi une majorité stricte ? Avec 7 pour, 5 contre et 5 abstentions, la
position la plus fréquente ne rassemble que $7/17\simeq41\%$. Les dix autres
votes ne sont pas dix infractions à une ligne : ils révèlent qu'aucune position
ne domine. Le scrutin est classé \emph{groupe partagé}, pas \emph{fracture}.

\subsection{Taux de dissidence}
\label{sec:dissidence}

Sur les seuls votes $\mathcal L_i$ où une ligne nette existe pour le groupe de
$i$, la dissidence est
\[
 \mathrm{dissidence}_i=
 \frac{\sum_{j\in\mathcal L_i}\ind(\text{vote de $i$}\neq
 \text{ligne de son groupe au scrutin $j$})}{|\mathcal L_i|}.
\]
Cette mesure décrit un écart à la majorité \emph{observée} du groupe : elle ne
prouve ni une opposition idéologique ni une désobéissance à une consigne privée.

\paragraph{Le seuil de comparabilité, et ce qu'il ne fait pas.} Une personne
avec trop peu de votes comparables est écartée du \emph{classement} — minimum de
50 par défaut, la valeur \code{analyze.MIN\_VOTES\_LIGNE}. La règle est
exactement celle de la présence (§\ref{sec:participation}) et elle a la même
portée : le taux reste calculé, reste publié, reste accompagné de son
dénominateur ; ce qu'il ne fait plus, c'est entrer dans la médiane, le maximum,
les rangs et la représentation graphique de la population.

Ce document a longtemps énoncé ce seuil pendant que le site n'en appliquait
aucun, et l'écart n'était pas théorique. Les trois seuls députés en exercice à
0~\% de dissidence de la 17\textsuperscript{e} législature sont précisément ceux
dont le groupe n'avait eu de ligne que sur 4, 8 et 25 scrutins : ils fixaient à
eux seuls l'extrémité « discipline » de la distribution, et chaque fiche
annonçait aux deux autres qu'ils étaient dépassés par des députés qui leur
étaient en réalité à égalité. Correctif porté au registre des corrections.

\subsubsection*{L'incertitude du taux}

Ce taux est une proportion estimée sur un nombre fini de votes : il a une marge,
et elle se calcule. Le projet publie l'intervalle de Jeffreys au niveau
$1-\gamma$, c'est-à-dire les quantiles de la loi
$\mathrm{Beta}(k+\tfrac12,\;|\mathcal L_i|-k+\tfrac12)$ où $k$ est le nombre de
votes dissidents.

Le choix de Jeffreys plutôt que de l'intervalle de Wald usuel
$\hat p\pm z\sqrt{\hat p(1-\hat p)/n}$ n'est pas indifférent ici. Wald se
dégrade précisément là où se trouve la majorité des députés : proche de zéro, il
produit des bornes négatives et son taux de couverture réel s'effondre. Or la
moitié de l'Assemblée est sous 3~\% de dissidence. L'intervalle de Jeffreys
reste dans $[0,1]$ par construction et se comporte correctement aux extrêmes.

\important{L'ordre de grandeur importe plus que la formule : l'intervalle fait
environ un point pour un député médian et jusqu'à quatre points pour les taux
élevés. Un rang de dissidence recouvre donc, en pratique, plusieurs dizaines de
positions que les données ne départagent pas. Le site publie le rang \emph{et}
le nombre de députés dont l'intervalle recouvre celui du député consulté.}

\prudence{\textbf{Cet intervalle suppose des votes indépendants, et ils ne le
sont pas.} La loi Beta ci-dessus traite les $|\mathcal L_i|$ votes d'un député
comme autant de tirages indépendants. Or les scrutins d'une même séance portent
sur le même texte dans le même rapport de force, et une rupture avec la ligne du
groupe se produit typiquement \emph{par épisode} — un député qui s'écarte sur un
texte s'écarte souvent sur la série d'amendements qui l'accompagne. Le nombre
d'observations réellement indépendantes est donc inférieur au nombre de votes,
et l'intervalle publié est \emph{plus étroit} que la variabilité réelle.

La conséquence est asymétrique et il faut la dire dans le bon sens : quand le
site annonce que $N$ députés ont un intervalle recouvrant celui du député
consulté, ce $N$ est un \textbf{minorant}. Un intervalle plus honnête en
recouvrirait davantage, jamais moins. Un recouvrement doit donc être lu comme
une preuve d'indiscernabilité, et une absence de recouvrement comme une
présomption faible — non l'inverse. Le bootstrap par blocs employé pour les
positions (§\ref{sec:bootstrap}) répondrait à cette objection ; l'appliquer à la
dissidence reste à faire.}

\subsection{Scrutins serrés et scrutins clivants}

Un scrutin est dit serré lorsque
\[
 e_j=|n_{\mathrm{pour},j}-n_{\mathrm{contre},j}|
\]
est faible (au plus 10 voix par défaut). À l'inverse, la part de votes dissidents
parmi les votes auxquels une ligne de groupe s'applique permet de construire un
taux de fracture du scrutin. Les deux notions sont différentes : un vote peut
être très serré mais parfaitement discipliné, ou largement adopté tout en
divisant fortement certains groupes.

\section{Deux cartes descriptives : ACP et positionnement multidimensionnel}

\subsection{Analyse en composantes principales (ACP)}
\label{sec:acp}

\subsubsection*{L'idée, avant la formule}

Chaque député est décrit par une liste de $m$ nombres : sa colonne de votes,
codée $+1$, $-1$ ou $0$. C'est un point dans un espace à $m$ dimensions, où $m$
vaut plusieurs milliers — donc invisible. L'ACP cherche le \emph{plan} sur
lequel projeter ce nuage en perdant le moins possible de sa forme, un peu comme
on cherche l'angle sous lequel photographier une sculpture pour qu'elle reste
reconnaissable.

Elle ne suppose rien sur la politique. Elle constate seulement que les colonnes
de votes ne se répartissent pas au hasard : si deux députés votent presque
toujours pareil, leurs points sont voisins, et le nuage s'étire dans certaines
directions plus que d'autres. L'ACP nomme ces directions.

\subsubsection*{Le calcul}

\textbf{Premier temps, centrer.} On retire à chaque colonne de $S$ sa moyenne :
\[
 (S_c)_{ij}=S_{ij}-\bar S_{\cdot j},
 \qquad \bar S_{\cdot j}=\frac1n\sum_{i}S_{ij}.
\]
Après cette opération, un coefficient ne dit plus « ce député a voté pour »
mais « ce député s'est écarté de la position moyenne du scrutin, et dans quel
sens ». C'est ce qu'on veut : un vote acquis à l'unanimité ne distingue
personne et devient une colonne de zéros, tandis qu'un vote serré garde toute
son amplitude.

\textbf{Second temps, décomposer.} On écrit la décomposition en valeurs
singulières
\[
 S_c=U\Sigma V^{\mathsf T},
\]
où $\Sigma$ est diagonale, de coefficients $\sigma_1\geq\sigma_2\geq\cdots\geq0$
appelés \emph{valeurs singulières}. Chaque colonne de $V$ est une direction de
l'espace des scrutins — une combinaison de votes ; la valeur singulière associée
mesure l'amplitude du nuage dans cette direction. Les coordonnées des députés
sur les deux premiers axes sont les deux premières colonnes de $U\Sigma$.

\textbf{L'inertie} est la dispersion totale du nuage, $\sum_q\sigma_q^2$. La
part portée par l'axe $r$ est
\[
 \frac{\sigma_r^2}{\sum_q\sigma_q^2}.
\]
Une part élevée sur le premier axe signifie qu'une seule direction résume
l'essentiel des écarts entre députés — c'est-à-dire qu'un clivage domine. Elle
ne dit rien de la \emph{nature} de ce clivage.

Les axes peuvent être retournés sans que le résultat change : $U\Sigma$ et
$-U\Sigma$ décrivent la même configuration. Leurs signes ne sont donc pas des
étiquettes politiques, et « à gauche du graphique » ne veut pas dire « à gauche
de l'hémicycle » autrement que par convention d'affichage.

\subsubsection*{La limite principale : ce que le zéro cache}

L'ACP travaille sur $S$, et $S$ code par $0$ quatre situations distinctes
(section~\ref{sec:signee}). Elle traite donc une absence \textbf{exactement
comme un point milieu} entre le pour et le contre. C'est une imputation, et
c'est la plus forte hypothèse de tout ce document.

Elle est d'autant plus lourde que les zéros sont majoritaires : sur la
17\textsuperscript{e} législature, environ trois quarts des cellules du cube
complet ne portent aucun suffrage exprimé. Les conséquences sont mesurables et
vont dans un sens prévisible :

\begin{itemize}
  \item un député souvent absent a une colonne proche de zéro ; après
  centrage, son point se retrouve près du centre du nuage, quelles que soient
  ses convictions ;
  \item la \emph{distance au centre} de la carte est donc corrélée à
  l'assiduité. Sur les données de référence, la corrélation entre la valeur
  absolue de la coordonnée sur l'axe 1 et le nombre de suffrages exprimés
  dépasse $0{,}7$ ;
  \item la zone centrale de la carte mêle en conséquence deux populations qu'il
  ne faut pas confondre : des députés réellement intermédiaires, et des députés
  peu présents.
\end{itemize}

Cette faiblesse est propre aux méthodes factorielles appliquées à des données
lacunaires ; elle n'affecte pas le taux d'accord de la
section~\ref{sec:proximite}, qui écarte explicitement les scrutins non partagés,
ni le modèle de points idéaux de la section~\ref{sec:ideal}, qui ne calcule la
vraisemblance que sur les votes observés.

\prudence{L'ACP répond correctement à « quelle est la structure d'ensemble des
votes ? » et incorrectement à « où se situe ce député ? ». Pour situer une
personne, c'est le modèle de points idéaux qu'il faut lire, avec son intervalle.
Le site n'emploie l'ACP que pour des vues d'ensemble, jamais sur une fiche
individuelle.}

\important{Une part d'inertie n'est pas une part de vérité. Elle est calculée
sur la matrice telle qu'elle a été codée, imputations comprises : elle mesure
la concentration du nuage, pas la qualité de la description.}

\subsection{Positionnement multidimensionnel (MDS)}

Le MDS répond à la même question par l'autre bout. Plutôt que de projeter des
colonnes de votes, il part directement des proximités deux à deux et cherche
une carte plane où les distances graphiques reproduisent au mieux les distances
mesurées. On pose $d_{ik}=1-T_{ik}$ : voter pareil, c'est être près.

L'intérêt est qu'il s'appuie sur $T_{ik}$, dont le dénominateur ne retient que
les scrutins réellement partagés — l'imputation massive qui affecte l'ACP ne
s'y produit donc pas. Une imputation subsiste néanmoins, mais bien plus
limitée : les paires sous le seuil de scrutins communs n'ont pas de taux, et
leur distance est provisoirement remplacée par la moyenne des distances
disponibles, faute de quoi la double centralisation serait impossible. Ces
paires sont peu nombreuses, et cette valeur de remplacement les place au centre
du nuage : la même prudence que pour l'ACP s'applique à elles, à cette échelle
réduite.

La matrice des distances est ensuite doublement centrée :
\[
 B=-\frac12 JD^{\circ2}J,\qquad
 J=I-\frac1n\mathbf{1}\mathbf{1}^{\mathsf T}.
\]
Les deux vecteurs propres associés aux plus grandes valeurs propres de $B$
donnent les coordonnées. Le MDS cherche une carte qui préserve autant que
possible les distances d'accord ; il est plus directement lié aux proximités
de paires, mais il est plus coûteux.

\prudence{Une carte rapproche des comportements observés. Elle ne démontre pas
une idéologie, une stratégie commune ou une cause du vote.}

\section{Le modèle de points idéaux : expliquer et prédire les votes binaires}
\label{sec:ideal}

\subsection{Le modèle statistique}

L'ACP décrit une configuration ; le modèle de points idéaux pose un mécanisme
probabiliste. Il ne conserve que les votes pour et contre. Pour un député $i$
et un scrutin $j$ :
\[
 Y_{ij}=\ind(\text{$i$ vote pour}),\qquad
 \Pp(Y_{ij}=1)=\operatorname{logit}^{-1}(\beta_j^{\mathsf T}x_i-\alpha_j),
\]
où
\[
 \operatorname{logit}^{-1}(z)=\frac{1}{1+e^{-z}}.
\]

\subsubsection*{Le vocabulaire, qui vient d'ailleurs}

Ce modèle est celui de la \emph{réponse à l'item}, construit à l'origine pour
analyser des examens : les élèves y ont une aptitude, les questions une
difficulté et un pouvoir discriminant. Le vocabulaire a suivi jusqu'ici, et
deux termes sur trois induisent en erreur si on les prend au pied de la lettre.

\begin{center}
\begin{tabular}{>{\raggedright\arraybackslash}p{.14\linewidth}p{.30\linewidth}p{.44\linewidth}}
\toprule
Symbole & Nom consacré & Ce que c'est ici \\
\midrule
$x_i\in\R^d$ & position \textbf{latente} & la position que l'on cherche. « Latente » signifie non observée : elle n'est écrite nulle part dans les données et n'est connue que par ses effets sur les votes. \\
$\beta_j\in\R^d$ & \textbf{discrimination} & la direction et la netteté avec lesquelles le scrutin $j$ sépare l'Assemblée. Rien à voir avec le sens ordinaire du mot. \\
$\alpha_j\in\R$ & \textbf{difficulté} & le terme le plus mal nommé : il ne mesure aucune difficulté, il situe le \emph{seuil} du vote — à partir de quelle position on bascule vers le « pour ». \\
\bottomrule
\end{tabular}
\end{center}

\subsubsection*{Comment lire l'équation}

Une image suffit à fixer les idées. Rangeons les députés sur une ligne. Chaque
scrutin y trace une frontière : à gauche on vote contre, à droite on vote pour.
Alors, en dimension 1 :

\begin{itemize}
  \item le \textbf{point de bascule} $\alpha_j/\beta_j$ (défini si
  $\beta_j\neq0$) dit \emph{où} passe la frontière — c'est la position à
  laquelle un député a exactement une chance sur deux de voter pour ;
  \item $|\beta_j|$ dit \emph{à quel point elle est nette}. Grand, la frontière
  est une coupure franche et la position prédit presque parfaitement le vote ;
  proche de zéro, le passage est un dégradé et le scrutin traverse les camps
  au lieu de les opposer.
\end{itemize}

Estimer le modèle revient donc à chercher le rangement des députés qui rend le
plus grand nombre de scrutins explicables par une telle frontière.

\important{Une position latente n'est pas une opinion : c'est la coordonnée qui
prédit le mieux les votes observés, sur les textes effectivement soumis pendant
la période. Changez les textes, vous changez l'axe.}

Les abstentions ne sont pas transformées en demi-votes. Elles sont exclues de
ce modèle parce qu'il est binaire, non parce qu'elles seraient dépourvues de
sens politique. Les analyses dédiées de la section~\ref{sec:abstention} les
traitent explicitement.

\subsection{Les filtres avant estimation}

Le modèle retire les scrutins avec trop peu de pour/contre, les députés avec
trop peu de votes binaires et, par défaut, les scrutins très peu contestés :
\[
 \frac{\min(n_{\mathrm{pour},j},n_{\mathrm{contre},j})}
 {n_{\mathrm{pour},j}+n_{\mathrm{contre},j}}\geq0{,}025.
\]
Un vote quasi unanime peut être politiquement important mais apporte très peu
d'information pour ordonner les députés sur un axe. Ce filtre sert à stabiliser
l'estimation, pas à déclarer ces scrutins sans intérêt.

\paragraph{L'assiette réelle du modèle, qui n'est pas celle des mesures
voisines.} Ces filtres ne sont pas marginaux, et il faut publier ce qu'ils
retirent. Sur la 17\textsuperscript{e} législature, le modèle est estimé sur
\textbf{127 des 245 votes qui engagent}. Le site a longtemps présenté la
position estimée comme calculée « sur les votes qui engagent », ce qui laissait
entendre 245 ; chaque fiche indique désormais le nombre de scrutins retenus et
le nombre de suffrages de ce député que le modèle a effectivement lus.

\textbf{Les 23 motions de censure en sortent toutes}, et ce n'est pas un effet
de réglage : la source n'y publie ni « contre » ni abstention — seuls les
soutiens de la motion votent (§\ref{sec:question}). La position minoritaire y
est donc nulle, la contestation vaut 0, et un scrutin auquel un seul camp prend
part ne peut pas séparer l'Assemblée le long d'un axe. C'est cohérent avec la
mise en garde déjà faite plus haut sur la lecture des motions de censure, et
c'est la raison pour laquelle elles comptent dans la présence sans compter dans
la position.

\subsection{Estimation alternée et pénalisation}

Le projet maximise une vraisemblance logistique pénalisée. Schématiquement, il
cherche à maximiser :
\[
 \ell=\sum_{i,j\ \text{observés}}
 \left[Y_{ij}\log(p_{ij})+(1-Y_{ij})\log(1-p_{ij})\right]
 -\frac{\lambda_x}{2}\sum_i\|x_i\|^2
 -\frac{\lambda_\beta}{2}\sum_j\|\beta_j\|^2.
\]
Les difficultés $\alpha_j$ ne sont pas pénalisées : un scrutin peut
légitimement être très largement acquis. Les paramètres sont ajustés par
régressions logistiques pénalisées alternées :

\begin{enumerate}
  \item positions $x_i$ fixées, estimer les paramètres de chaque scrutin ;
  \item paramètres des scrutins fixés, estimer les positions des députés ;
  \item recommencer jusqu'à stabilisation de la log-vraisemblance.
\end{enumerate}

Chaque régression logistique est mise à jour par IRLS (\emph{iteratively
reweighted least squares}), c'est-à-dire un pas de Newton utilisant les poids
$p_{ij}(1-p_{ij})$. Le code traite les nombreux scrutins en parallèle dans des
produits matriciels, ce qui rend possible le bootstrap.

La pénalisation de $\beta_j$ est importante : avec une séparation parfaite,
le maximum de vraisemblance non pénalisé tend vers $|\beta_j|=\infty$. On
obtiendrait des frontières verticales et des positions artificiellement tassées.

\paragraph{Comment $\lambda_\beta$ est choisi, et pourquoi ce n'est pas le
maximum d'APRE.} La valeur retenue est $\lambda_\beta=0{,}5$. Elle n'est pas
celle qui maximise l'ajustement hors échantillon, et il serait malhonnête de
laisser croire le contraire : la validation croisée donne, sur les données de
référence,

\begin{center}
\begin{tabular}{ccc}
\toprule
$\lambda_\beta$ & APRE hors échantillon & $\max_j|\beta_j|$ \\
\midrule
0,02 & 0,825 & 26,9 \\
0,10 & 0,824 & 13,5 \\
\textbf{0,50} & \textbf{0,812} & \textbf{7,7} \\
2,00 & 0,807 & 5,0 \\
8,00 & 0,776 & 2,9 \\
\bottomrule
\end{tabular}
\end{center}

La règle de choix est donc explicite, et elle a deux termes : \emph{la plus
faible pénalisation dont les discriminations restent bornées}. À
$\lambda_\beta=0{,}02$, on gagne 1,3 point d'APRE et $|\beta_j|$ monte à 27 —
c'est-à-dire des frontières quasi verticales, une séparation numériquement
presque parfaite sur certains scrutins, et des intervalles de bootstrap dont la
largeur devient dictée par quelques votes. Nous préférons un demi-point
d'ajustement à des positions dont la stabilité ne survivrait pas au
rééchantillonnage. C'est un arbitrage, il est discutable, et c'est précisément
pourquoi le tableau est publié plutôt que la seule conclusion.

\subsection{Le problème d'identification}

Le produit $\beta_j^{\mathsf T}x_i$ reste le même si l'on change simultanément
l'échelle et l'orientation de l'espace. Après estimation, le projet :

\begin{enumerate}
  \item centre les positions ;
  \item les réduit à un écart-type unitaire ;
  \item les tourne sur leurs axes principaux ;
  \item oriente le premier axe pour qu'un groupe d'ancrage soit du côté négatif.
\end{enumerate}

L'ancrage rend deux exécutions lisibles dans le même sens. Il ne donne pas une
signification intrinsèque (« gauche », « droite ») à l'axe : il fixe seulement
le miroir.

\subsection{Évaluer l'ajustement : classification et APRE}

La classification compare $\ind(p_{ij}\geq0{,}5)$ au vote observé. Mais elle
favorise les scrutins très déséquilibrés : prédire toujours la majorité est déjà
souvent correct. L'APRE (\emph{aggregate proportional reduction in error})
compare donc les erreurs du modèle aux erreurs de ce prédicteur naïf :
\[
 \mathrm{APRE}=\frac{E_0-E_{\text{modèle}}}{E_0},
\]
où $E_0$ est la somme des effectifs minoritaires de chaque scrutin, et
$E_{\text{modèle}}$ le nombre d'erreurs de classification. APRE $=0$ signifie
« pas mieux que toujours choisir la majorité » ; APRE $=1$ signifie aucune erreur.

Pour choisir le nombre de dimensions et la pénalisation, une part des cellules
de vote est mise de côté au hasard. Le modèle est estimé sur le reste puis
évalué sur ces votes non vus. Une dimension supplémentaire qui n'améliore que
l'ajustement d'apprentissage est du \emph{surajustement} — elle a mémorisé du
bruit propre à l'échantillon — et ne doit pas être lue comme un nouveau clivage
politique.

Le retrait porte sur des \textbf{cellules} et non sur des scrutins entiers, et
ce choix est contraint : retirer un scrutin entier priverait le modèle de ses
paramètres $\alpha_j$ et $\beta_j$, si bien qu'il n'y aurait plus rien à
prédire dessus.

\prudence{Il faut reconnaître une tension entre cette procédure et le bootstrap
de la section~\ref{sec:bootstrap}. Ce dernier tire des blocs entiers au motif —
exact — que les scrutins d'une même séance ne sont pas indépendants. La
validation croisée, elle, retire des cellules isolées : le modèle voit alors la
plupart des votes du même scrutin et peut en estimer précisément les
paramètres, ce qui rend la prédiction des cellules retirées plus facile qu'une
véritable prévision hors échantillon. Le protocole est celui d'usage pour ce
type de modèle, mais il est plus indulgent que le bootstrap sur la même donnée.
La conclusion sur le nombre de dimensions doit être lue comme une borne
supérieure de ce qui est justifié, non comme une démonstration.}

\subsection{Combien de dimensions, et combien le site en publie}
\label{sec:dimensions}

Ces deux questions n'ont pas la même réponse, et les confondre a produit sur le
site une affirmation fausse qu'il faut redresser ici.

\textbf{Ce que le test dit.} Sur les données de référence, l'APRE hors
échantillon vaut 0,81 à une dimension et 0,93 à deux. Le gain est large,
reproductible d'un tirage à l'autre, et très supérieur à l'écart entre
répétitions. Une seconde dimension \emph{n'est pas} du surajustement : elle
prédit mieux des votes que le modèle n'a pas vus. Substantiellement, le premier
axe oppose LFI au reste et le second le RN au reste — deux oppositions
systématiques qui ne portent pas sur les mêmes textes et ne peuvent donc pas
tenir sur une seule ligne. À trois dimensions l'APRE monte encore (0,94), mais
l'écart entre apprentissage et test se creuse, et nous nous arrêtons à deux.

\textbf{Ce que le site publie.} Une seule dimension. Le site l'a longtemps
justifié en écrivant qu'une seconde serait « décorative » — c'est-à-dire en
présentant comme un résultat de test l'exact contraire de ce que le test donne.
La justification correcte est autre, et elle est éditoriale : tout ce qu'une
fiche affirme d'une position se lit sur une ligne — un rang, un intervalle, un
nombre de voisins indiscernables, une bande. Deux axes affichés ensemble
forment un plan, c'est-à-dire une carte que le lecteur lira comme une carte des
opinions, ce que ce projet refuse de publier (§\ref{sec:acp}).

\prudence{La conséquence doit être connue du lecteur, et le site la publie
désormais : la position affichée sur une fiche résume par une seule coordonnée
une Assemblée qui en demande au moins deux. Deux députés que le second axe
sépare nettement peuvent apparaître au même endroit du premier. La position
publiée n'est pas fausse — elle est incomplète d'une dimension dont nous savons
qu'elle existe.}

\subsection{Incertitude par bootstrap de blocs}

\label{sec:bootstrap}
Une position estimée est une estimation, non une position mesurée sans erreur.
Le projet tire avec remise des \emph{blocs} de scrutins, réestime le modèle à
chaque tirage, puis prend les quantiles empiriques. Pour un niveau $1-\gamma$,
l'intervalle est :
\[
 [\widehat x_i^{(\gamma/2)},\widehat x_i^{(1-\gamma/2)}].
\]

Les scrutins d'une même séance sont corrélés : ils peuvent porter sur le même
texte et le même rapport de force. Les tirer indépendamment reviendrait à
surestimer la quantité d'information. Le bootstrap regroupe donc, si possible,
par dossier législatif, sinon par séance, sinon par date. Les intervalles sont
plus larges que sous un bootstrap scrutin par scrutin, mais plus honnêtes.

\important{Deux intervalles qui se recouvrent ne démontrent pas l'égalité des
positions, mais ils doivent empêcher de lire un petit écart de rang comme une
différence certaine.}

\section{L'abstention est une position à analyser}\label{sec:abstention}

\subsection{Taux d'abstention}

Pour un groupe, un député, une catégorie ou une portée, le taux est
\[
 \frac{\#\{\text{abstentions}\}}
 {\#\{\text{pour, contre ou abstention}\}}.
\]
Les absences et non-votes ne sont pas dans le dénominateur : sinon on mesurerait
un mélange d'abstention volontaire et d'absentéisme.

\subsection{Consigne, retrait ou indétermination}

Chaque abstention est jointe à la ligne recalculée de son groupe. Elle est
classée :

\begin{itemize}
  \item \textbf{consigne}, si le groupe a une ligne nette d'abstention ;
  \item \textbf{retrait}, si le groupe a une ligne nette pour ou contre ;
  \item \textbf{indéterminée}, si aucune ligne nette n'existe ou si le groupe a
  trop peu voté.
\end{itemize}

Les parts sont divisées par le \emph{total de toutes les abstentions}, y compris
les indéterminées. Cette troisième catégorie mesure une limite de classement ;
la retirer du dénominateur ferait artificiellement gonfler les deux premières.

\subsection{Tester l'idée d'une position intermédiaire}

Le modèle de points idéaux est estimé sans abstentions. Pour chaque scrutin où
les trois populations sont assez nombreuses, on compare les médianes des
positions : $m_+$ pour les pour, $m_-$ pour les contre et $m_0$ pour les
abstentionnistes. Après avoir posé $b=\min(m_+,m_-)$ et $h=\max(m_+,m_-)$, la
position relative est
\[
 r=\frac{m_0-b}{h-b}.
\]
Une valeur entre 0 et 1 place la médiane des abstentionnistes entre les deux
camps ; $r=0{,}5$ la place à mi-chemin. C'est un test \emph{hors modèle} :
l'information d'abstention n'a pas servi à construire l'axe. Il ne prouve pas
que chaque abstention est un compromis ; il évalue cette proposition scrutin
par scrutin.

\subsection{Quand l'abstention aurait pu faire basculer l'issue}

Si un texte est adopté, $e=n_{\mathrm{pour}}-n_{\mathrm{contre}}>0$. Il suffit
de déplacer $e$ abstentionnistes vers le contre pour atteindre l'égalité, qui
fait échouer l'adoption : le seuil est donc $a\geq e$.

Si le texte est rejeté, $e=n_{\mathrm{contre}}-n_{\mathrm{pour}}\geq0$. Le pour
doit dépasser le contre, pas seulement l'égaler : le seuil est $a>e$, soit au
moins $e+1$ abstentionnistes. La fonction ne prétend pas connaître leurs
intentions ; elle établit seulement qu'ils étaient assez nombreux pour modifier
l'issue sous ce scénario arithmétique.

\section{Amendements et réseau de cosignatures}
\label{sec:cosign}

\subsection{Volumes et adoption}

Pour un député ou un groupe, le projet compte les amendements dont l'auteur est
identifié. La source publie pour chacun un \emph{sort} — adopté, rejeté, tombé,
retiré, non soutenu — ou rien du tout lorsque le texte n'a pas encore été
discuté. En notant $A$ l'ensemble des amendements déposés par l'auteur et
$A_{\mathrm{ex}}\subseteq A$ ceux dont la source publie un sort :
\[
 \mathrm{taux\ d'adoption}=\frac{\#\{a\in A_{\mathrm{ex}} : \text{sort}(a)=\text{adopté}\}}
 {\#A_{\mathrm{ex}}},
 \qquad \text{indéfini si } A_{\mathrm{ex}}=\varnothing .
\]

\textbf{Le dénominateur est l'examen, pas le dépôt}, et l'écart n'est pas
marginal : sur la 17\textsuperscript{e} législature, 54~485 des 123~224
amendements enregistrés n'ont pas de sort publié, soit 44~\%. Diviser par les
dépôts revient à compter comme des échecs des amendements dont l'Assemblée n'a
rien dit ; le taux d'un député baisserait alors à chaque nouveau dépôt sur un
texte en cours de navette, sans qu'aucun vote ait eu lieu. Le nombre
d'amendements en attente est publié à côté du taux, pour que la différence
entre $\#A$ et $\#A_{\mathrm{ex}}$ reste lisible.

Ce taux est descriptif : l'adoption dépend notamment de la recevabilité, de la
procédure, de la majorité, du sujet et du statut de l'auteur. Comparer seulement
des taux sans le volume de dépôts serait trompeur.

\subsection{Construire les liens}

Un amendement signé par A, B et C crée les trois liens non orientés A--B,
A--C et B--C. Si $D$ est la matrice d'incidence députés $\times$ amendements,
avec $D_{ia}=1$ si le député $i$ a signé l'amendement $a$, alors
\[
 W=DD^{\mathsf T}.
\]
Pour $i\neq k$, $W_{ik}$ est le nombre d'amendements cosignés ensemble. Sur la
diagonale, $W_{ii}$ est le nombre total d'amendements signés par $i$ : ce n'est
pas un lien avec soi-même.

Le réseau exclut par défaut les amendements de rapporteur et les amendements
ayant plus de dix signataires. Le premier filtre évite de confondre une tâche
institutionnelle avec une alliance. Le second évite qu'un dépôt collectif de
groupe crée mécaniquement tous les liens internes possibles.

\subsection{Affinité de Jaccard}

Le compte $W_{ik}$ favorise les députés les plus actifs. L'affinité utilisée
pour classer les paires est donc l'indice de Jaccard :
\[
 J_{ik}=\frac{W_{ik}}{s_i+s_k-W_{ik}},
\]
où $s_i=W_{ii}$ est le nombre d'amendements signés par $i$. Il s'agit de la
taille de l'intersection des deux répertoires de signatures divisée par leur
union. Les députés avec moins de 20 signatures sont masqués : une petite base
peut produire un Jaccard de 1 sans signal robuste.

\subsection{Part des liens entre groupes}

Pour deux groupes distincts $G$ et $H$, le volume de liens est
\[
 L_{GH}=\sum_{i\in G}\sum_{k\in H}W_{ik}.
\]
Pour un lien interne, la matrice symétrique le compte deux fois et la diagonale
doit être retirée :
\[
 L_{GG}=\frac{1}{2}\left(\sum_{i\in G}\sum_{k\in G}W_{ik}
 -\sum_{i\in G}W_{ii}\right).
\]
La part de la production de liens de $G$ orientée vers $H$ est
\[
 \pi_{G\to H}=\frac{L_{GH}}{\sum_Q L_{GQ}}.
\]
La même convention -- chaque paire une fois, jamais la diagonale -- est
employée au numérateur et au dénominateur. Ainsi, les parts d'une ligne totalisent
1 et les groupes deviennent comparables.

\subsection{Les « courtiers » : sortir de son groupe plus que ne l'impose sa taille}

La part brute de liens hors groupe est mécaniquement plus grande dans un petit
groupe. Pour le député $i$ du groupe $G$, le projet la compare à l'attendu sous
mélange aléatoire :
\[
 \text{part observée}_i=\frac{\text{liens de $i$ hors de $G$}}
 {\text{tous les liens de $i$}},\qquad
 \mathrm{part\ attendue}_i=\frac{N-|G|}{N-1}.
\]
Le ratio observé/attendu vaut 1 sous cette référence. Un ratio supérieur à 1
signale une ouverture hors groupe plus forte que ce que la seule taille du groupe
produirait. Les non-inscrits sont exclus de cette mesure : « hors de son groupe »
n'a alors pas de référence bien définie.

\section{Détecter les sujets qui montent}\label{sec:sujets}

\subsection{Corpus et unités de comptage}

Le corpus hebdomadaire contient les titres de scrutins et, lorsque disponible,
les titres et les premiers caractères des exposés sommaires d'amendements. Les
textes sont passés en minuscules sans accents, tokenisés en mots et bigrammes,
puis filtrés : mots vides, jargon légistique, termes de procédure et noms de
députés ne doivent pas devenir des « sujets ».

L'unité pertinente est le \textbf{document qui mentionne au moins une fois un
terme}, non le nombre d'occurrences. Si un mot est répété dix fois dans un long
document, ce document contribue une fois à $n_{t,w}^{\mathrm{docs}}$. Avec
$N_t$ documents durant la semaine $t$, la fréquence comparable est
\[
 f_{t,w}=\frac{n_{t,w}^{\mathrm{docs}}}{N_t}.
\]
Cela empêche qu'une semaine comportant des textes longs ou beaucoup de dépôts
soit confondue avec un changement de thème.

\subsection{Attendu binomial et score de poussée}

Sur les semaines de référence $R$, la fréquence de référence du terme $w$ est
\[
 p_w=\frac{\sum_{t\in R}n_{t,w}^{\mathrm{docs}}}{\sum_{t\in R}N_t}.
\]
Dans la semaine courante de taille $N$, l'effectif attendu est
$\mu_w=Np_w$. Sous l'hypothèse de stabilité, le nombre de documents mentionnant
le terme suit approximativement
\[
 X_w\sim\mathrm{Binomiale}(N,p_w).
\]
Le projet calcule la probabilité de queue $\Pp(X_w\geq x_w)$ et un score de
classement
\[
 z_w=\frac{x_w-\mu_w}{\sqrt{\mu_w+1}}.
\]
Le score classe l'ampleur de la hausse ; la $p$-valeur vérifie si elle est
compatible avec le hasard sous le modèle binomial. Pour un terme absent de la
référence, une petite fréquence plancher évite une $p$-valeur artificiellement
nulle sur une seule mention.

\subsection{Tests multiples : Benjamini--Hochberg}

Des milliers de mots et d'expressions peuvent être testés chaque semaine. Au
seuil de 5\% sans correction, on attendrait beaucoup de faux signaux. Les
$p$-valeurs ordonnées $p_{(1)}\leq\cdots\leq p_{(m)}$ sont converties en
$q$-valeurs de Benjamini--Hochberg :
\[
 q_{(r)}=\min_{s\geq r}\left(\frac{m}{s}p_{(s)}\right),
\]
bornée par 1. Filtrer à $q\leq0{,}05$ contrôle le taux de fausses découvertes
attendu à 5\% parmi les signaux retenus, sous les hypothèses usuelles de la
procédure. Enfin, des n-grammes parfaitement redondants sont fusionnés pour ne
pas faire passer plusieurs fragments d'un même titre pour plusieurs sujets.

\section{Lire les sorties sans leur faire dire plus qu'elles ne disent}\label{sec:lire}

\begin{center}
\begin{tabular}{>{\raggedright\arraybackslash}p{.25\linewidth}p{.31\linewidth}p{.31\linewidth}}
\toprule
Indicateur & Ce qu'il décrit & Ce qu'il ne démontre pas \\
\midrule
Accord de vote & Concordance de positions exprimées & Motif commun, alliance durable, identité idéologique. \\
Participation & Vote aux scrutins où l'élu pouvait voter & Volume total de travail parlementaire. \\
Dissidence & Écart à une majorité observée de groupe & Violation d'une consigne privée ou désaccord idéologique. \\
Point idéal & Axe latent qui prédit des votes binaires & Étiquette intrinsèque ou position politique totale. \\
Cosignature & Travail d'initiative rendu visible par les signatures & Alliance générale ou causalité entre députés. \\
Sujet montant & Hausse relative de mentions dans le corpus & Importance normative ou attention de toute la société. \\
\bottomrule
\end{tabular}
\end{center}

Une bonne lecture procède dans cet ordre :

\begin{enumerate}
  \item vérifier le périmètre (dates, portée, élus inclus) ;
  \item lire le numérateur et le dénominateur ;
  \item comparer, si nécessaire, les scrutins de détail et les votes qui engagent ;
  \item regarder les seuils d'exclusion et l'incertitude ;
  \item revenir au titre, au texte et au contexte politique du scrutin.
\end{enumerate}

\section{Reproduire et auditer}\label{sec:reproduire}

Les calculs décrits ici sont implémentés dans :

\begin{itemize}
  \item \code{src/radar/parse.py} : normalisation, catégories, portées et
  positions de groupe ;
  \item \code{src/radar/analyze.py} : cube, accord, cohésion, participation,
  dissidence et cartes ;
  \item \code{src/radar/ideal.py} : modèle de points idéaux, validation et
  bootstrap ;
  \item \code{src/radar/abstention.py} : analyses spécifiques aux abstentions ;
  \item \code{src/radar/cosign.py} : réseau de cosignatures ;
  \item \code{src/radar/topics.py} : corpus et détection de poussées.
\end{itemize}

Pour reconstruire les tables depuis les données ouvertes de l'Assemblée :

\begin{verbatim}
uv run radar update
\end{verbatim}

Les tests unitaires du projet contiennent également des cubes synthétiques :
ils vérifient notamment les faux 100\% d'accord, la différence entre moyenne de
taux et quotient de sommes, l'orientation du modèle et l'APRE.

\important{La reproductibilité ne consiste pas seulement à pouvoir relancer le
code. Elle suppose de publier les définitions, les seuils, les exclusions et
les incertitudes qui rendent un chiffre discutable.}

\section{Limites connues}

Les limites qui suivent ne sont pas des réserves de style. Chacune borne une
conclusion que le lecteur pourrait être tenté de tirer.

\textbf{Le scrutin public n'est pas le vote.} L'immense majorité des décisions
de l'Assemblée se prend à main levée, sans trace nominative. Les scrutins
publics sont un sous-ensemble non aléatoire : ils sont demandés, souvent
précisément parce qu'un groupe veut rendre des positions visibles. Tout ce qui
est mesuré ici porte sur cette fraction, et sur elle seule.

\textbf{La participation mesurée n'est pas la présence.} Un député présent en
commission, en circonscription ou en mission n'apparaît pas dans nos
dénominateurs. Le taux publié dit à quelle fraction des scrutins publics
auxquels il pouvait prendre part un élu a exprimé une position ; il ne dit
rien de son travail parlementaire pris dans son ensemble.

\textbf{Le point idéal est un artefact utile, pas une opinion.} L'axe estimé
est la direction qui prédit le mieux les votes observés. Son orientation est
fixée par convention, son échelle est arbitraire, et son interprétation
substantielle dépend entièrement des textes soumis au vote pendant la période.
Deux législatures ne sont pas comparables sur cette échelle.

\textbf{L'appartenance de groupe n'est datée qu'à moitié.} La table des votes
porte le groupe du député \emph{au moment du scrutin}, et c'est ce groupe-là qui
sert à reconstruire la ligne et à compter la dissidence : ces deux mesures sont
exactes historiquement. En revanche, les agrégats qui passent par la table des
députés — accord entre groupes, cohésion, cosignatures — rattachent un vote
passé au groupe \emph{actuel} de son auteur. Sur les données de référence, cela
concerne 19\,033 votes sur 1\,270\,476, soit 1,5~\%, répartis sur 348 députés.
Une version antérieure de ce document affirmait que les changements de groupe
étaient partout traités à leur date d'effet ; c'était vrai de la dissidence et
faux des agrégats, et la formulation est corrigée ici plutôt qu'atténuée.

\textbf{La délégation n'est pas neutralisée dans la dissidence.} La part
déléguée est désormais définie et publiée (section~\ref{sec:delegation}), mais
un seul indicateur en tire les conséquences : la présence, qui l'affiche à côté
de son numérateur. Le taux de dissidence, lui, traite un vote délégué comme
tous les autres, alors qu'un écart à la ligne du groupe a pu être matériellement
produit par le collègue porteur. Recalculé sur les seuls votes émis en personne,
le classement déplace le député médian de quelques rangs, et bien davantage ceux
qui délèguent souvent. Publier les deux versions, ou au moins cette sensibilité,
reste à faire.

\textbf{La participation reste publiée sans intervalle.} La position latente est
accompagnée de son intervalle de bootstrap (section~\ref{sec:bootstrap}) et la
dissidence du sien (section~\ref{sec:dissidence}) ; le taux de présence, lui,
est encore un point nu. C'est aussi une proportion, son intervalle se calcule de
la même façon, et rien ne justifie l'asymétrie sinon qu'elle n'a pas encore été
levée. La lacune est plus étroite qu'elle n'y paraît — le dénominateur vaut 245
pour les trois quarts des députés, ce qui donne des intervalles serrés — mais
elle est réelle pour les élus arrivés en cours de législature.

\textbf{La cosignature n'est pas la collaboration.} Signer un amendement est un
acte peu coûteux, parfois collectif par défaut au sein d'un groupe. Le réseau
décrit une pratique de signature, dont le rapport au travail réel est une
question ouverte que ces données ne tranchent pas.

\section{Disponibilité des données, du code et de ce document}

\textbf{Données sources.} Jeux de données ouverts de l'Assemblée nationale,
\url{https://data.assemblee-nationale.fr}, sous Licence Ouverte
\textsc{Etalab}. Le projet enregistre pour chaque archive l'en-tête
\code{Last-Modified} renvoyé par le serveur ; la date de la source est publiée
sur le site à côté de la date de calcul.

\textbf{Code.} Le paquet \code{radar} est publié sous licence
\textsc{agpl}-3.0-or-later. Le choix de cette licence est délibéré : un
outil dont l'objet est la transparence ne peut pas être réutilisé pour
produire des chiffres dont la méthode resterait fermée.

\textbf{Tables intermédiaires.} Elles ne sont pas encore archivées : chaque
mise à jour reconstruit les tables et remplace les précédentes. Reconstruire
l'état du site à une date passée suppose donc aujourd'hui que les archives
amont soient encore disponibles dans la même version — ce que rien ne
garantit. L'archivage d'un instantané par mise à jour est prévu ; tant qu'il
n'est pas en place, cette limite est réelle et il faut la connaître.

\textbf{Ce document.} Il est publié au format \textsc{pdf} avec le site, à
l'adresse \url{https://radar-parlementaire.fr/methode/note-methodologique.pdf},
et sa source \LaTeX{} est versionnée avec le code.

\section{Déclaration sur l'usage d'outils d'assistance à la rédaction}
\label{sec:ia}

Ce document, ainsi qu'une part substantielle du code du projet, ont été
rédigés avec l'assistance de \textbf{Claude}, modèle de langage développé par
\textbf{Anthropic}, utilisé de manière interactive.

Le périmètre de cette assistance a été large : formulation et structuration du
texte, mise en forme \LaTeX{}, formalisation mathématique des définitions,
rédaction et révision du code Python, et travail critique sur les choix de
méthode — plusieurs incohérences entre indicateurs, dont deux définitions
divergentes d'une même grandeur, ont été identifiées au cours de ces échanges.

La responsabilité du contenu n'en est pas partagée pour autant. Les choix de
méthode, les seuils, les décisions de ne pas publier une mesure jugée trop
fragile, la vérification des résultats et l'ensemble du texte publié relèvent
de l'auteur, qui en répond seul. Une erreur dans ce document est une erreur de
l'auteur.

Cette déclaration figure ici pour la même raison que tout le reste du
document. Un projet qui exige des autres qu'ils exposent leur méthode ne peut
pas dissimuler la sienne, y compris lorsqu'elle porte sur ses propres outils
de travail. Mentionner un tel usage devient une pratique attendue dans
l'édition scientifique ; nous considérons qu'elle l'est tout autant dès lors
qu'on publie des chiffres sur la vie politique.

\section*{Licence de ce document}
\addcontentsline{toc}{section}{Licence de ce document}

Ce texte est mis à disposition sous licence
\href{https://creativecommons.org/licenses/by/4.0/deed.fr}
     {Creative Commons Attribution 4.0 International} (\textsc{cc by} 4.0).
Il peut être cité, reproduit, traduit et adapté, y compris pour le contredire,
à condition d'en attribuer la paternité.

\begin{thebibliography}{9}
\addcontentsline{toc}{section}{Références}

\bibitem{benjamini1995}
Y.~Benjamini et Y.~Hochberg.
\newblock Controlling the False Discovery Rate: A Practical and Powerful
  Approach to Multiple Testing.
\newblock \emph{Journal of the Royal Statistical Society, Series B},
  57(1):289--300, 1995.

\bibitem{clinton2004}
J.~Clinton, S.~Jackman et D.~Rivers.
\newblock The Statistical Analysis of Roll Call Data.
\newblock \emph{American Political Science Review}, 98(2):355--370, 2004.

\bibitem{hix2006}
S.~Hix, A.~Noury et G.~Roland.
\newblock Dimensions of Politics in the European Parliament.
\newblock \emph{American Journal of Political Science}, 50(2):494--511, 2006.

\bibitem{jaccard1901}
P.~Jaccard.
\newblock Étude comparative de la distribution florale dans une portion des
  Alpes et du Jura.
\newblock \emph{Bulletin de la Société vaudoise des sciences naturelles},
  37:547--579, 1901.

\bibitem{kunsch1989}
H.~R.~Künsch.
\newblock The Jackknife and the Bootstrap for General Stationary Observations.
\newblock \emph{The Annals of Statistics}, 17(3):1217--1241, 1989.

\bibitem{poole1985}
K.~T.~Poole et H.~Rosenthal.
\newblock A Spatial Model for Legislative Roll Call Analysis.
\newblock \emph{American Journal of Political Science}, 29(2):357--384, 1985.

\bibitem{poole1997}
K.~T.~Poole et H.~Rosenthal.
\newblock \emph{Congress: A Political-Economic History of Roll Call Voting}.
\newblock Oxford University Press, 1997.

\bibitem{rice1928}
S.~A.~Rice.
\newblock \emph{Quantitative Methods in Politics}.
\newblock Alfred A. Knopf, New York, 1928.

\bibitem{torgerson1952}
W.~S.~Torgerson.
\newblock Multidimensional Scaling: I. Theory and Method.
\newblock \emph{Psychometrika}, 17(4):401--419, 1952.

\end{thebibliography}

\end{document}

